\chapter{Le buy side}\label{ch:m1:the-buy-side}

Quatre cents dollars quittent le bulletin de paie d'un enseignant le dernier jour du mois
et arrivent dans son fonds de pension. Le conseil d'administration du fonds a décidé que soixante
pour cent de ses capitaux suivent un indice mondial d'actions, qu'un quart est confié à
des gérants rémunérés pour battre celui-ci, et qu'un vingtième soit placé dans des \omterm{def:m1:what-a-trading-firm-does:hedge-fund}{hedge funds}. Au
moment où ses quatre cents dollars sont investis, ils ont été répartis entre une
machine qui doit acheter tout ce que détient l'indice, un sélectionneur de titres qui est
évalué chaque trimestre, et une société en commandite qui garde un cinquième de ce qu'elle
gagne. Aucun des trois n'a décidé de négocier aujourd'hui en raison d'une opinion sur les prix.
Ce chapitre porte sur les personnes situées de l'autre côté du téléphone du dealer~:
qui elles sont, ce pour quoi elles sont rémunérées, et --- la question
qui intéresse un trader --- \emph{pourquoi elles négocient}.

\section{Propriétaires d'actifs et sociétés de gestion}

\begin{definition}[Buy side]\label{def:m1:the-buy-side:buy-side}
Le \emph{buy side}\index{buy side} est l'ensemble des institutions qui investissent
de l'argent en vue d'un rendement et achètent des services de trading pour ce faire~: \emph{propriétaires d'actifs}
(fonds de pension, assureurs, fonds souverains, fonds de dotation, ménages), qui supportent
le résultat final, et les \emph{\omterm{def:m1:what-a-trading-firm-does:asset-manager}{sociétés de gestion}} et les fonds qu'ils engagent.
\end{definition}

\begin{definition}[Mandat]\label{def:m1:the-buy-side:mandate}
Un \emph{mandat}\index{mandat} est le contrat par lequel un propriétaire d'actifs
confie un portefeuille à un gérant. Il fixe l'univers investissable, le
\omterm{def:m1:the-buy-side:benchmark}{indice de référence}, les limites de risque, la commission et ce que le gérant
n'est pas autorisé à faire (effet de levier, vente à découvert, produits dérivés, concentration).
\end{definition}

\begin{remark}[Les fonds publics sont encadrés par la loi]\label{rem:m1:the-buy-side:law}
Un fonds vendu au grand public porte un \omterm{def:m1:the-buy-side:mandate}{mandat} rédigé par le
législateur. Un fonds américain qui se présente comme \emph{diversifié} au titre du
Investment Company Act of 1940 doit investir 75~\% de ses actifs en veillant à ce
qu'aucun émetteur ne dépasse 5~\% du fonds ni 10~\% des titres assortis de droit de vote de cet émetteur. Un
fonds UCITS européen peut placer au plus 10~\% dans un même émetteur, et ses positions
supérieures à 5~\% ne peuvent pas totaliser plus de 40~\%. Ces règles, et non des opinions
sur les prix expliquent une bonne part de ce que détiennent les plus grands fonds ---
et pourquoi ils doivent vendre un titre dont le cours a simplement trop monté.
\end{remark}

\begin{omfigure}
\begin{tikzpicture}[font=\small,
  box/.style={draw=omDef,rounded corners=2pt,minimum width=2.4cm,minimum height=0.9cm,align=center,fill=omDef!4},
  flow/.style={-{Stealth[length=2mm]},thick,omDef}]
\node[box] (sv) at (0,0) {Épargnant};
\node[box] (ow) at (3.5,0) {Propriétaire d'actifs};
\node[box] (mg) at (7.0,0) {Société de gestion\\ou fonds};
\node[box] (sl) at (10.5,0) {Sell side};
\node[box] (mk) at (13.5,0) {Marché};
\draw[flow] (sv) -- (ow);
\draw[flow] (ow) -- node[above=11pt,font=\footnotesize]{mandat} (mg);
\draw[flow] (mg) -- node[above=11pt,font=\footnotesize]{ordres} (sl);
\draw[flow] (sl) -- (mk);
\node[omThm,font=\footnotesize,align=center] at (3.5,-1.05) {frais\\administratifs};
\node[omThm,font=\footnotesize,align=center] at (7.0,-1.05) {frais de gestion\\et de performance};
\node[omThm,font=\footnotesize,align=center] at (10.5,-1.05) {commissions, spread,\\financement};
\end{tikzpicture}

\omcaption{La chaîne qui va de l'épargnant au marché. Chaque maillon est rémunéré (en rouge), et
chacun ajoute une raison de négocier qui n'a rien à voir avec le cours
de n'importe quel titre.}
\end{omfigure}

\section{Indices de référence et tracking error}

\begin{definition}[Indice de référence, rendement actif, tracking error]\label{def:m1:the-buy-side:benchmark}
Un \emph{indice de référence}\index{indice de référence} est le portefeuille de référence, en général un
indice publié, par rapport auquel un \omterm{def:m1:the-buy-side:mandate}{mandat} est mesuré. Si $r^P_t$ et $r^B_t$
sont les rendements du portefeuille et de l'indice, le \emph{rendement actif}
est $r^A_t = r^P_t - r^B_t$, et la
\emph{tracking error}\index{tracking error} est l'écart-type annualisé
de $r^A$.
\end{definition}

\begin{definition}[Gestion passive et gestion active]\label{def:m1:the-buy-side:passive}
La \emph{gestion passive}\index{gestion passive} vise une \omterm{def:m1:the-buy-side:benchmark}{tracking error}
nulle~: le gérant détient l'\omterm{def:m1:the-buy-side:benchmark}{indice de référence} et perçoit quelques points de base pour
le faire à moindre coût. La \emph{gestion active}\index{gestion active} accepte
une \omterm{def:m1:the-buy-side:benchmark}{tracking error} dans l'espoir d'un rendement actif moyen positif.
\end{definition}

\begin{proposition}[Tracking error à partir des pondérations actives]\label{prop:m1:the-buy-side:te}
Soient $w^P, w^B$ les vecteurs de pondérations du portefeuille et de l'\omterm{def:m1:the-buy-side:benchmark}{indice de référence}, et $\Sigma$
la matrice de covariance annualisée des rendements des actifs. Avec les pondérations actives
$a = w^P - w^B$ (dont la somme est nulle), la \omterm{def:m1:the-buy-side:benchmark}{tracking error} est
$\mathrm{TE} = \sqrt{a^\top \Sigma\, a}$.
\end{proposition}
\begin{proof}
$r^A = (w^P - w^B)^\top r = a^\top r$, dont la variance est $a^\top\Sigma a$.
\end{proof}

\begin{example}[Vingt paris en surpondération]\label{ex:m1:the-buy-side:bets}
Un gérant surpondère 20 actions de 1~\% chacune et sous-pondère 20 autres
de 1~\% chacune. Si les rendements propres à chaque titre sont indépendants, avec une volatilité de 25~\%
et que les facteurs communs s'annulent, $\mathrm{TE} = \sqrt{40 \times (0.01 \times
0.25)^2} = 1.6\%$. Un \omterm{def:m1:the-buy-side:mandate}{mandat} qui plafonne la \omterm{def:m1:the-buy-side:benchmark}{tracking error} à 3~\% laisse
peu de marge au-delà de cela~: un portefeuille «~actif~» n'est, pour l'essentiel, que l'indice.
\end{example}

\begin{definition}[Ratio d'information]\label{def:m1:the-buy-side:ir}
Le \emph{ratio d'information}\index{ratio d'information} d'un \omterm{def:m1:the-buy-side:mandate}{mandat} est le
rendement actif annuel moyen divisé par la \omterm{def:m1:the-buy-side:benchmark}{tracking error}~:
$\mathrm{IR} = \E[r^A]/\mathrm{TE}$.
\end{definition}

\begin{proposition}[Combien de temps il faut pour établir un talent de gestion]\label{prop:m1:the-buy-side:years}
Avec des rendements actifs indépendants et identiquement distribués, et un \omterm{def:m1:the-buy-side:ir}{ratio
d'information} véritable $\mathrm{IR}$, la statistique $t$ du rendement actif moyen
après $T$ années est, en espérance, $\mathrm{IR}\sqrt{T}$. Atteindre $t = 2$
exige $T = 4/\mathrm{IR}^2$ années.
\end{proposition}
\begin{proof}
Sur $T$ années, la moyenne empirique du rendement actif annuel a une
erreur-type $\mathrm{TE}/\sqrt{T}$, donc $t = \E[r^A]\sqrt{T}/\mathrm{TE} =
\mathrm{IR}\sqrt{T}$.
\end{proof}

\begin{example}[Seize ans]\label{ex:m1:the-buy-side:sixteen}
Un gestionnaire long only dont le $\mathrm{IR}$ est réellement bon, égal à 0.5, a besoin de 16
ans pour que son historique se distingue du hasard au seuil usuel~; avec $\mathrm{IR} = 0.25$, il lui en faut 64.
Une stratégie de \omterm{def:m1:what-a-trading-firm-does:market-maker}{teneur de marché} dont le ratio de Sharpe quotidien équivaut à $\mathrm{IR} = 8$
requiert environ trois semaines. Cette seule phrase explique pourquoi les gestionnaires de long terme
sont embauchés sur des récits et renvoyés sur du bruit, tandis que les sociétés à horizon court peuvent
être pilotées à partir de la statistique
(\cref{fig:m1:the-buy-side:years}).
\end{example}

\begin{omfigure}
\begin{tikzpicture}
\begin{axis}[width=10.5cm,height=4.8cm,ymode=log,
  xlabel={ratio d'information},ylabel={années pour atteindre $t=2$},
  tick label style={font=\small},label style={font=\small},
  xmin=0.2,xmax=3,ymin=0.4,ymax=110,grid=major,grid style={gray!25},
  ytick={0.5,1,2,5,10,20,50,100},yticklabels={0.5,1,2,5,10,20,50,100},
  table/col sep=comma]
\addplot[omDef,very thick] table[x=ir,y=years]{figdata/markets-1/03-the-buy-side/years_needed.csv};
\addplot[only marks,mark=*,omThm] coordinates {(0.5,16) (1,4) (2,1)};
\node[omThm,font=\footnotesize,anchor=west] at (axis cs:0.55,19) {16 ans};
\node[omThm,font=\footnotesize,anchor=west] at (axis cs:1.05,5) {4 ans};
\node[omThm,font=\footnotesize,anchor=west] at (axis cs:2.05,1.25) {1 an};
\end{axis}
\end{tikzpicture}

\omcaption{Historique nécessaire pour établir la compétence, $T = 4/\mathrm{IR}^2$
(axe vertical logarithmique). Données~: calculées par le
script du chapitre.}\label{fig:m1:the-buy-side:years}
\end{omfigure}

\begin{dated}{2026-09}{Dans quelle mesure le marché est devenu passif}\label{dat:m1:the-buy-side:passive}
À la fin de 2025, les fonds communs indiciels et les fonds indiciels cotés (ETF) détenaient
52~\% des actifs des fonds américains à long terme, une majorité pour la deuxième année, selon
le décompte de l'association professionnelle du secteur des fonds.
\end{dated}

\section{Hedge funds et leurs frais}

Un \omterm{def:m1:what-a-trading-firm-does:hedge-fund}{hedge fund} (\cref{def:m1:what-a-trading-firm-does:hedge-fund}) échappe aux
règles des fonds publics parce qu'il n'est proposé qu'à des investisseurs
professionnels et fortunés~; il peut emprunter, vendre à découvert et utiliser
des produits dérivés, et sa rémunération est organisée autrement.

\begin{definition}[Frais de gestion, commission de performance, high-water mark]\label{def:m1:the-buy-side:fees}
Les \emph{frais de gestion}\index{frais de gestion} sont un pourcentage annuel fixe de
les actifs. Une \emph{commission de performance}\index{commission de performance} est un pourcentage
du bénéfice de la période. Avec un
\emph{high-water mark}\index{high-water mark}, la commission de performance n'est prélevée
que sur la part dont la valeur liquidative de l'investisseur dépasse son
plus haut niveau antérieur de paiement~: les pertes doivent être comblées avant que le
gestionnaire ne soit de nouveau rémunéré sur les gains.
\end{definition}

\begin{definition}[Commodity trading advisor]\label{def:m1:the-buy-side:cta}
Un \emph{commodity trading advisor}\index{commodity trading advisor} (CTA) est,
en droit américain, une société enregistrée pour conseiller sur le trading de contrats à terme~; dans l'usage du secteur,
un fonds qui négocie des contrats à terme de façon systématique sur l'ensemble des classes d'actifs, le plus souvent
en suivant les tendances (livre~8).
\end{definition}

\begin{method}[Calcul des frais avec un high-water mark]\label{met:m1:the-buy-side:hwm}
On part d'une valeur liquidative $N_0 = 1$ et d'un niveau de référence $H_0 = 1$. Chaque année, avec
un rendement brut $g_t$, un taux de \omterm{def:m1:the-buy-side:fees}{frais de gestion} $\mu$ et un taux de \omterm{def:m1:the-buy-side:fees}{commission de performance} $\phi$~:
\begin{enumerate}
\item prélever les \omterm{def:m1:the-buy-side:fees}{frais de gestion} $\mu N_{t-1}$~;
\item calculer $N' = N_{t-1}(1+g_t) - \mu N_{t-1}$~;
\item prélever la \omterm{def:m1:the-buy-side:fees}{commission de performance} $\phi\,\max(0,\, N' - H_{t-1})$~;
\item fixer $N_t = N'$ diminué de cette commission, et $H_t = \max(H_{t-1}, N_t)$.
\end{enumerate}
Les conventions varient (frais calculés sur les actifs d'ouverture ou moyens,
cristallisation trimestrielle, taux de rendement minimal)~: il faut consulter le document de souscription.
\end{method}

\begin{remark}[Une commission de performance est une option]\label{rem:m1:the-buy-side:option}
Le gain $\phi\max(0, N'-H)$ est celui de $\phi$ options d'achat sur la valeur du fonds
de prix d'exercice égal au \omterm{def:m1:the-buy-side:fees}{high-water mark}, accordées chaque année au gestionnaire
sans contrepartie. Une option vaut d'autant plus que la volatilité est forte~: un gestionnaire sans
aucune compétence perçoit des frais espérés plus élevés en prenant davantage de risque
(\cref{fig:m1:the-buy-side:convex}), et un gestionnaire très en dessous de son \omterm{def:m1:the-buy-side:fees}{high-water mark} détient
une option si hors de la monnaie que fermer le fonds et en lancer un autre est le choix rationnel.
Les investisseurs répondent par des limites de risque, par des clawbacks,
et en exigeant que le gestionnaire place son propre argent dans le fonds.
\end{remark}

\begin{omfigure}
\begin{tikzpicture}
\begin{axis}[width=10.5cm,height=4.6cm,
  xlabel={volatilité annuelle du fonds (\%)},ylabel={frais moyens (\% par an)},
  tick label style={font=\small},label style={font=\small},
  xmin=0,xmax=30,ymin=1.5,ymax=3,grid=major,grid style={gray!25},table/col sep=comma]
\addplot[omDef,very thick,mark=*,mark size=1.3pt] table[x=vol_pct,y=fee_pct_per_year]{figdata/markets-1/03-the-buy-side/fee_vs_vol.csv};
\end{axis}
\end{tikzpicture}

\omcaption{Frais totaux d'un fonds «~2 and 20~» avec \omterm{def:m1:the-buy-side:fees}{high-water mark} et
un rendement brut espéré \emph{nul}, en fonction de la volatilité qu'il supporte. La compétence est
absente par construction~: la hausse représente la valeur de l'option. Données~:
simulation réalisée par le script du chapitre, 4\,000 trajectoires de dix ans par
point.}\label{fig:m1:the-buy-side:convex}
\end{omfigure}

\begin{dated}{2026-09}{Taille du secteur des hedge funds}\label{dat:m1:the-buy-side:hf}
Les \omterm{def:m1:what-a-trading-firm-does:hedge-fund}{hedge funds} géraient environ 5.6~billions de dollars à la fin du deuxième trimestre de
2026, après avoir franchi pour la première fois 5~billions de dollars à la fin de 2025,
selon la base de données sectorielle la plus citée. Cela représente moins de 4~\%
des actifs gérés professionnellement (\cref{dat:m1:what-a-trading-firm-does:aum}),
mais une part bien plus grande des volumes échangés, parce que le capital est mis
à effet de levier et tourne plusieurs fois par an.
\end{dated}

\section{Pourquoi chacun d'eux négocie}

La première question d'un \omterm{def:m1:what-a-trading-firm-does:market-maker}{teneur de marché} face à un ordre est de savoir si l'expéditeur
sait quelque chose (\cref{def:m1:what-a-trading-firm-does:adverse-selection}). Les
motivations du \omterm{def:m1:the-buy-side:buy-side}{buy side} se répartissent le long de cette ligne.

\begin{example}[Motifs, et ce qu'ils révèlent à la contrepartie]\label{ex:m1:the-buy-side:motives}
\begin{center}\small
\begin{tabular}{@{}p{3.3cm}p{5.3cm}p{5.6cm}@{}}
\hline
Qui & Pourquoi l'ordre existe & Ce qu'il véhicule \\
\hline
Fonds indiciel & souscriptions ou rachats~; l'indice a changé & aucune opinion~; taille et date prévisibles \\
Fonds de pension, assureur & rééquilibrage vers les poids cibles~; couverture des passifs & aucune opinion sur le titre~; vend ce qui a monté \\
Long only actif & une opinion sur quelques mois & une information partielle, lente à peser \\
\omterm{def:m1:what-a-trading-firm-does:hedge-fund}{Hedge fund}, approche fondamentale & une opinion sur quelques semaines~; un événement & information~; urgence autour des annonces \\
\omterm{def:m1:what-a-trading-firm-does:hedge-fund}{Hedge fund}, approche systématique & un signal sur quelques heures à quelques jours & information éphémère~; nombreux petits ordres \\
Tout fonds à effet de levier & un appel de marge ou une limite de risque & aucune opinion --- mais forcé, de grande taille et corrélé à celui des autres \\
Clientèle de détail & épargne, opinions, divertissement & presque rien \\
\hline
\end{tabular}
\end{center}
La troisième colonne explique pourquoi un dealer traite différemment des ordres identiques
venant de clients différents, et pourquoi les courtiers vendent le flux de détail
(\cref{ch:m1:retail-flow-and-wholesaling}).
\end{example}

\begin{remark}[Un flux prévisible est la stratégie de quelqu'un d'autre]\label{rem:m1:the-buy-side:predictable}
Les flux non informés \emph{et} prévisibles --- changements d'indice, rééquilibrage de fin de mois,
remise à zéro quotidienne des fonds à effet de levier --- sont la matière première des
stratégies événementielles du livre~8. Les \omterm{def:m1:the-buy-side:mandate}{mandats} du \omterm{def:m1:the-buy-side:buy-side}{buy side} sont
publics~; on peut dire, en première approximation, qu'il en va de même de ses transactions.
\end{remark}

\section{Tutoriel~: trois fonds, dix ans}\label{tut:m1:the-buy-side:threefunds}

\begin{tutorial}
\textbf{Objectif.} Mesurer le \omterm{def:m1:the-buy-side:benchmark}{tracking error} sur des rendements simulés, puis calculer
les frais d'un \omterm{def:m1:what-a-trading-firm-does:hedge-fund}{hedge fund} année après année. \textbf{État final~:} le graphique ci-dessous et le
tableau du problème du week-end.
\begin{tutsteps}
\item \textbf{\omterm{def:m1:the-buy-side:benchmark}{Tracking error} et \omterm{def:m1:the-buy-side:ir}{ratio d'information}.}
\omcode{code/markets-1/03-the-buy-side/python/buyside.py}{8}{21}{Tracking error, ratio d'information et nombre d'années nécessaires pour établir la compétence.}\label{lst:m1:the-buy-side:te}
\item \textbf{Contrôle sur des fonds simulés.} Le script de la figure trace dix ans
de rendements mensuels de l'\omterm{def:m1:the-buy-side:benchmark}{indice de référence}, un fonds indiciel (5 points de base de frais) et
un fonds actif construit avec 1~\% d'alpha brut, des frais de 0.75~\% et
un \omterm{def:m1:the-buy-side:benchmark}{tracking error} de 4~\%. Mesurés~: des \omterm{def:m1:the-buy-side:benchmark}{tracking errors} de 0.11~\% et 3.78~\%, et
un rendement réalisé de $+0.66\%$ par an pour le fonds actif, contre un
rendement vrai de $+0.25\%$ après frais. Dix ans ne disent presque rien sur lequel est
lequel --- comme le prédit la \cref{prop:m1:the-buy-side:years}.
\item \textbf{Frais avec \omterm{def:m1:the-buy-side:fees}{high-water mark}.}
\omcode{code/markets-1/03-the-buy-side/python/buyside.py}{32}{51}{La méthode de ce chapitre, ligne à ligne.}\label{lst:m1:the-buy-side:hwm}
\item \textbf{Calculer la trajectoire sur dix ans} du problème du week-end. Le fonds perçoit des \omterm{def:m1:the-buy-side:fees}{commissions de performance} dans
trois années seulement sur dix~; les sept autres sont passées
sous la \omterm{def:m1:the-buy-side:fees}{high-water mark}.
\end{tutsteps}
\textbf{Pour aller plus loin.} Porter les \omterm{def:m1:the-buy-side:fees}{frais de gestion} à 1~\% et la \omterm{def:m1:the-buy-side:fees}{commission de
performance} à 30~\%~: qui y gagne~? Ajouter ensuite un seuil de rendement minimal~: aucune
\omterm{def:m1:the-buy-side:fees}{commission de performance} sur les premiers 4~\% de rendement annuel.
\end{tutorial}

\begin{omfigure}
\begin{tikzpicture}
\begin{axis}[width=11cm,height=5.2cm,
  xlabel={année},ylabel={valeur d'un placement de 1},
  tick label style={font=\small},label style={font=\small},
  xmin=0,xmax=10,ymin=0.8,ymax=2.0,grid=major,grid style={gray!25},
  legend columns=3,legend style={font=\footnotesize,at={(0.5,-0.32)},anchor=north,draw=none,fill=none,/tikz/every even column/.append style={column sep=8pt}},
  table/col sep=comma]
\addplot[omExo,thick,dashed] table[x=year,y=gross]{figdata/markets-1/03-the-buy-side/hwm_path.csv};
\addplot[omDef,very thick] table[x=year,y=net]{figdata/markets-1/03-the-buy-side/hwm_path.csv};
\addplot[omThm,thick,const plot] table[x=year,y=hwm]{figdata/markets-1/03-the-buy-side/hwm_path.csv};
\legend{avant frais,valeur nette de l'investisseur,high-water mark}
\end{axis}
\end{tikzpicture}

\omcaption{Dix ans d'un fonds «~2 et 20~». La \omterm{def:m1:the-buy-side:fees}{high-water mark} ne monte que lorsque la valeur nette
atteint un nouveau plus haut~; l'écart entre la courbe en tirets et la courbe pleine est
ce que coûtent à l'investisseur les frais, et la croissance qu'ils auraient rapportée.
Données~: la trajectoire de rendements du problème du week-end, calculée par le script du chapitre.}
\end{omfigure}

\section{Construction~: le moteur des frais}\label{bld:m1:the-buy-side:fees}

\begin{build}
\bfield{Rôle} La société miniature gérera un fonds. Elle doit calculer
ce que doivent ses investisseurs, par investisseur et par période. \bfield{Interface} \texttt{FeeTerms(mgmt, perf, hurdle=0.0,
periods\_per\_year=1)} et \texttt{accrue(terms, state, gross\_return)}
renvoyant un nouvel \texttt{InvestorState(nav, high\_water)} et les deux frais.
\bfield{Règles} Un état par investisseur et par date de souscription~: deux investisseurs
d'un même fonds ont des \omterm{def:m1:the-buy-side:fees}{high-water marks} distinctes. Les \omterm{def:m1:the-buy-side:fees}{frais de gestion} sont prélevés sur
la valeur d'ouverture, au prorata de chaque période. Un seuil de rendement relève la \omterm{def:m1:the-buy-side:fees}{high-water mark} du
taux de seuil à chaque période, avant le calcul de la \omterm{def:m1:the-buy-side:fees}{commission de performance}. Les frais ne sont jamais
négatifs.
\bfield{Tests d'acceptation} \texttt{code/firm/fees/tests/}~: reproduit la
trajectoire du tutoriel~; aucune \omterm{def:m1:the-buy-side:fees}{commission de performance} sous la \omterm{def:m1:the-buy-side:fees}{high-water mark}~; avec des
périodes trimestrielles et volatilité nulle, la commission annuelle égale celle de la règle annuelle.
\bfield{Pour aller plus loin} Enregistrer les frais dans le grand livre du
\cref{ch:m1:what-a-trading-firm-does}, en produits de la \omterm{def:m1:what-a-trading-firm-does:asset-manager}{société de gestion}.
\end{build}

\begin{omsources}
\item Investment Company Institute, \emph{2026 Investment Company Fact Book},
chapitre 2.
\item HFR, \emph{Global Hedge Fund Industry Report}, deuxième trimestre 2026, et
les commentaires de marché qui l'accompagnent.
\item Services de la US Securities and Exchange Commission, \emph{Report to Congress
regarding threshold limits applicable to diversified companies}, février
2022 (le test 75--5--10 de l'Investment Company Act of 1940).
\item Directive 2009/65/EC (UCITS), article 52, et ESMA, \emph{Questions and
answers on the application of the UCITS Directive}.
\item R.~Grinold et R.~Kahn, \emph{Active Portfolio Management}, 2e éd.,
McGraw-Hill, 2000, chapitres 4 et 5 (ratio d'information, tracking error).
\item W.~Goetzmann, J.~Ingersoll et S.~Ross, «~High-water marks and hedge
fund management contracts~», \emph{Journal of Finance} 58 (2003).
\end{omsources}

\section{Exercices}

\begin{exercise}[$\star$]\label{exo:m1:the-buy-side:1}
Un fonds a rapporté 9.1~\%, 3.4~\%, $-6.2\%$, 12.0~\% en quatre ans alors que son
\omterm{def:m1:the-buy-side:benchmark}{indice de référence} a rendu 8.0~\%, 4.0~\%, $-7.0\%$, 10.5~\%. Calculer les
rendements actifs, leur moyenne, le \omterm{def:m1:the-buy-side:benchmark}{tracking error} (écart-type empirique) et le
\omterm{def:m1:the-buy-side:ir}{ratio d'information}.
\end{exercise}

\begin{exercise}[$\star$]\label{exo:m1:the-buy-side:2}
Combien d'années d'historique faut-il pour atteindre $t = 2$ avec un
\omterm{def:m1:the-buy-side:ir}{ratio d'information} de 0.4~? de 1.5~? Et combien de \emph{jours de négociation} avec un ratio annualisé
de 6~?
\end{exercise}

\begin{exercise}[$\star$]\label{exo:m1:the-buy-side:3}
Un fonds indiciel de 50~milliards de dollars prélève 4~points de base~; un \omterm{def:m1:what-a-trading-firm-does:hedge-fund}{hedge fund} de 2~milliards de dollars
prélève 2 et 20 et rapporte 10~\% brut. Quel gérant gagne le plus cette année~?
\end{exercise}

\begin{exercise}[$\star\star$]\label{exo:m1:the-buy-side:4}
Un fonds OPCVM (UCITS) détient une action à 9.5~\%, trois à 6~\% chacune et deux à 5.5~\%
chacune~; toute autre position est inférieure à 5~\%. Le fonds respecte-t-il la règle 5/10/40~?
L'action à 9.5~\% progresse de 30~\% alors que le reste du fonds reste stable~: quel est son
nouveau poids, et que doit faire le gérant~?
\end{exercise}

\begin{exercise}[$\star\star$]\label{exo:m1:the-buy-side:5}
Un gérant gère 50 surpondérations et 50 sous-pondérations de 0.5~\% chacune~;
la volatilité spécifique aux titres est de 30~\% et les rendements spécifiques sont indépendants.
Calculer le \omterm{def:m1:the-buy-side:benchmark}{tracking error}. Il vise 3~\%~: quel poids actif par position
permet de l'atteindre~?
\end{exercise}

\begin{exercise}[$\star\star$]\label{exo:m1:the-buy-side:6}
Un fonds sous contrat «~2 et 20~» avec un \omterm{def:m1:the-buy-side:fees}{high-water mark} rapporte
$+30\%$, $-20\%$, $+10\%$ brut. Appliquer la \cref{met:m1:the-buy-side:hwm} et
donner, pour chaque année, les deux frais (gestion et performance) et la valeur nette de clôture.
\end{exercise}

\begin{exercise}[$\star\star\star$]\label{exo:m1:the-buy-side:7}
\emph{Programmation.} Avec \texttt{run\_fees}, comparer «~2 et 20~» avec «~1 et
30~» sur le scénario de dix ans du problème du week-end. Indiquer les frais totaux sous chaque
contrat et la valeur finale de l'investisseur. Quel contrat préfère un investisseur qui
croit en son gérant, et pourquoi le gérant pourrait-il préférer l'autre~?
\end{exercise}

\begin{exercise}[$\star\star\star$]\label{exo:m1:the-buy-side:8}
\emph{Trouver la faille.} Un consultant classe 400 fonds actifs selon leur rendement
actif sur cinq ans, retient les 20 meilleurs et affirme que «~les gérants retenus
ont un \omterm{def:m1:the-buy-side:ir}{ratio d'information} démontré supérieur à 1~». Le fonds type de
l'échantillon a un \omterm{def:m1:the-buy-side:benchmark}{tracking error} de 4~\%. Supposer qu'\emph{aucun} gérant ne possède
une quelconque compétence. Estimer le rendement actif du vingtième meilleur fonds et son
\omterm{def:m1:the-buy-side:ir}{ratio d'information} apparent, et énoncer la faille.
\end{exercise}

\section{Problème~: Deux et vingt}

\begin{problem}[{Problème du week-end --- dix ans dans un hedge fund}]\label{pb:m1:the-buy-side:1}
Un fonds de dotation investit 100~millions de dollars dans un \omterm{def:m1:what-a-trading-firm-does:hedge-fund}{hedge fund} prélevant des \omterm{def:m1:the-buy-side:fees}{frais de
gestion} de 2~\% sur la valeur d'ouverture et une \omterm{def:m1:the-buy-side:fees}{commission de performance} de 20~\% avec un \omterm{def:m1:the-buy-side:fees}{high-water mark}
(la \cref{met:m1:the-buy-side:hwm}). Sur dix ans, les rendements bruts du fonds
sont
\[
+20~\%,\ -15~\%,\ +10~\%,\ +25~\%,\ -5~\%,\ +8~\%,\ +30~\%,\ -20~\%,\ +15~\%,\ +12~\%.
\]

\medskip\noindent\textbf{Partie I --- Les trois premières années.}
\begin{enumerate}
\item Année 1~: calculer les \omterm{def:m1:the-buy-side:fees}{frais de gestion}, la \omterm{def:m1:the-buy-side:fees}{commission de performance}, la valeur
nette de clôture et le nouveau \omterm{def:m1:the-buy-side:fees}{high-water mark}.
\item Année 2~: mêmes calculs. Pourquoi n'y a-t-il aucune \omterm{def:m1:the-buy-side:fees}{commission de performance}~?
\item Année 3~: le fonds gagne 10~\% brut. Une \omterm{def:m1:the-buy-side:fees}{commission de performance} est-elle due~?
\item Après trois ans, combien l'investisseur a-t-il gagné, et combien le
gérant a-t-il perçu~?
\end{enumerate}

\medskip\noindent\textbf{Partie II --- Le parcours complet} (utiliser le code du tutoriel
ou un tableur).
\begin{enumerate}[resume]
\item Au cours de quelles années une \omterm{def:m1:the-buy-side:fees}{commission de performance} est-elle versée~?
\item Donner le total des \omterm{def:m1:the-buy-side:fees}{frais de gestion} et le total des \omterm{def:m1:the-buy-side:fees}{commissions de performance} sur
les dix ans, en millions de dollars.
\item Donner la valeur finale de l'investisseur et le rendement net annualisé.
\item Donner la valeur finale brute de 100~millions de dollars et le rendement
brut annualisé.
\item Les frais versés et le profit de l'investisseur ne font pas le total du
profit brut. Calculer l'écart et l'expliquer.
\end{enumerate}

\medskip\noindent\textbf{Partie III --- Qui a obtenu quoi.}
\begin{enumerate}[resume]
\item Quelle fraction du profit brut a été versée au gérant~? Quelle
fraction l'investisseur a-t-il conservée~?
\item Un fonds indiciel facturant 5~points de base a rendu 6~\% par an brut sur
la même décennie. Comparer le résultat de l'investisseur.
\item Les rendements du \omterm{def:m1:what-a-trading-firm-does:hedge-fund}{hedge fund} étaient décorrélés des actions. Donner une
raison pour laquelle le fonds de dotation pourrait néanmoins être satisfait.
\item L'écart-type des rendements bruts annuels du fonds est de 17~\%.
Estimer son ratio de Sharpe brut en prenant la moyenne des rendements bruts et un
taux sans risque nul. Combien d'années faudrait-il pour l'établir à $t = 2$~?
\end{enumerate}

\medskip\noindent\textbf{Partie IV --- Incitations.}
\begin{enumerate}[resume]
\item À la fin de l'année 8, la valeur nette est très inférieure au \omterm{def:m1:the-buy-side:fees}{high-water mark}. De combien
le fonds doit-il progresser, net des \omterm{def:m1:the-buy-side:fees}{frais de gestion}, avant que le gérant
touche à nouveau une \omterm{def:m1:the-buy-side:fees}{commission de performance}~?
\item Expliquer pourquoi le gérant pourrait alors (a) accroître le risque ou (b) fermer le
fonds, et qui perd dans chaque cas.
\item Un second investisseur souscrit au début de l'année 9. Quel est
\emph{son} \omterm{def:m1:the-buy-side:fees}{high-water mark}, et quelle \omterm{def:m1:the-buy-side:fees}{commission de performance} paie-t-il au cours de l'année 9~?
\item Pourquoi cela rend-il inévitable la comptabilité par investisseur
(la \cref{bld:m1:the-buy-side:fees})~?
\item En utilisant la \cref{fig:m1:the-buy-side:convex}, estimer les frais annuels perçus
par un gérant dépourvu de compétence à une volatilité de 20~\%, et l'excédent
par rapport aux seuls \omterm{def:m1:the-buy-side:fees}{frais de gestion}.
\item Énoncer le \emph{résultat nommé}~: la part du profit brut
revenant au gérant sur les dix ans, en pourcentage.
\item En deux phrases, dire ce qu'un investisseur doit négocier en premier~: les
2, les 20, ou autre chose.
\end{enumerate}
\end{problem}

\section{Questions d'entretien}

\begin{interviewq}[$\star$ \iqroles{trader, chercheur, banque}]\label{iq:m1:the-buy-side:1}
Qu'est-ce que la \omterm{def:m1:the-buy-side:benchmark}{tracking error}, et pourquoi un gérant convaincu de la qualité d'une action refuserait-il
de la laisser représenter 15~\% du portefeuille~?
\end{interviewq}

\begin{interviewq}[$\star$ \iqroles{chercheur, ingénieur ML}]\label{iq:m1:the-buy-side:2}
Une stratégie affiche un ratio de Sharpe de 1. Quelle durée de backtest vous faut-il
pour croire que ce ratio n'est pas nul~?
\end{interviewq}

\begin{interviewq}[$\star\star$ \iqroles{trader, chercheur}]\label{iq:m1:the-buy-side:3}
Vous recevez deux ordres identiques de vente de 500\,000 actions, l'un émanant
d'un fonds indiciel et l'autre d'un \omterm{def:m1:what-a-trading-firm-does:hedge-fund}{hedge fund} fondamental. Les pricez-vous de la même manière~?
\end{interviewq}

\begin{interviewq}[$\star\star$ \iqroles{chercheur, banque}]\label{iq:m1:the-buy-side:4}
Expliquez le \omterm{def:m1:the-buy-side:fees}{high-water mark} à un non-spécialiste, puis expliquez ce qu'il fait
à l'appétit du gérant pour le risque.
\end{interviewq}

\begin{interviewq}[$\star\star$ \iqroles{chercheur, trader}]\label{iq:m1:the-buy-side:5}
Les fonds indiciels détiennent plus de la moitié des actifs des fonds américains. Nommez deux
flux de transactions prévisibles que cela engendre, et qui pourrait tirer profit de chacun.
\end{interviewq}

\begin{interviewq}[$\star\star\star$ \iqroles{chercheur, ingénieur ML}]\label{iq:m1:the-buy-side:6}
Sur 1\,000 stratégies sans aucune compétence, quel est le meilleur résultat qu'on puisse observer
après trois ans~? Donnez un ordre de grandeur de son ratio de Sharpe.
\end{interviewq}
